\documentclass[11pt,a4paper]{amsart}
\usepackage{amsmath,amssymb}
\usepackage[margin=30mm]{geometry}
\usepackage[hidelinks]{hyperref}

\hypersetup{
  pdftitle={The One-Step Ascent Set of Jeffery's Collatz Map},
  pdfauthor={Yunkun Chen}
}

\newtheorem{theorem}{Theorem}[section]
\newtheorem{lemma}[theorem]{Lemma}
\newcommand{\Npos}{\mathbb{Z}_{>0}}

\title[The ascent set of Jeffery's Collatz map]
  {The One-Step Ascent Set of Jeffery's Collatz Map}
\author{Yunkun Chen}
\date{November 25, 2026}
\subjclass[2020]{Primary 11B83. Secondary 11D61, 68V20.}
\keywords{Collatz map, Beatty sequence, powers of two and three, formal proof, Lean}

\begin{document}

\begin{abstract}
L.~Edson Jeffery conjectured that the array defining OEIS A254312
enumerates, without repetition, the positive integers \(m\) for which his
Collatz-derived map \(S\) satisfies \(S(m)>m\). We prove this conjecture.
The row containing an ascent \(m\) is determined by
\(1+\nu_2(3m-2)\), and a congruence modulo \(3\) determines its column.
An explicit separation estimate for powers of \(2\) and \(3\) gives the
strict ascent criterion needed in the proof. A Lean~4 formalization
verifies the array bijection and the upward-antidiagonal enumeration.
\end{abstract}

\maketitle

\section{Introduction}

For \(z\in\Npos\), let \(\nu_2(z)\) be the exponent of \(2\) in \(z\).
The accelerated odd Collatz map is
\[
  F(x)=\frac{3x+1}{2^{\nu_2(3x+1)}}\qquad (x\in\Npos,\ x\text{ odd}).
\]
Jeffery's map, as defined in OEIS A257480 \cite{jeffery-map}, is
\begin{equation}\label{eq:step}
  S(m)=\frac{3+(3/2)^{\nu_2(1+F(4m-3))}(1+F(4m-3))}{6}
  \qquad(m\in\Npos).
\end{equation}
OEIS A254311 records the positive integers \(m\) with \(S(m)>m\)
\cite{jeffery-ascents}. For the Beatty sequence A117630 \cite{beatty},
write
\[
 a(n)=\left\lfloor\frac{n\log 3}{\log(3/2)}\right\rfloor.
\]
Jeffery proposed that the array in A254312 \cite{jeffery-array}, with
entries
\begin{equation}\label{eq:array}
 A(n,k)=\frac{1}{6}
 \left[
  2^{a(n)}\left(6k-\frac{(3-(-1)^{a(n)})(1-(-1)^n)}{2}\right)
  -2^n+4
 \right],\qquad n,k\in\Npos,
\end{equation}
enumerates the same set without repetition.

\begin{theorem}\label{thm:main}
The map \((n,k)\mapsto A(n,k)\) is a bijection from
\(\Npos^2\) onto \(\{m\in\Npos:S(m)>m\}\). If \(A(n,k)=m\), then
\[
 n=1+\nu_2(3m-2).
\]
In particular, A254312 is a permutation of A254311.
\end{theorem}

The proof uses a one-sided gap estimate for \(2^{n+v}-3^v\).
Kayadibi \cite[Theorem~6.2]{kayadibi} proves a stronger estimate for
\(v\ge2\) at the least admissible exponent \(n+v\).
We include a short derivation from Ellison's separation theorem
\cite[Theorem~3]{ellison}. Appendix~\ref{sec:formal-gap} gives the
derivation used in the Lean development, which takes its
large-exponent bound from Merle \cite{merle}.

\newpage
\section{A gap estimate and the ascent criterion}

\begin{lemma}\label{lem:gap}
If \(n,v\in\Npos\) and \(3^v<2^{n+v}\), then
\begin{equation}\label{eq:gap}
  2^{n+v}-3^v\ge 2^n-1.
\end{equation}
\end{lemma}

\begin{proof}
For fixed \(v\), the difference between the two sides of \eqref{eq:gap},
\(2^n(2^v-1)-3^v+1\), increases with \(n\). It suffices to take
\(n=x-v\), where \(x\) is the least integer with \(2^x>3^v\).
The inequalities \(2^v<3^v<2^{2v}\) give \(v<x\le 2v\).
For \(1\le v\le17\), the exact integer calculations in
Table~\ref{tab:small} prove \eqref{eq:gap}. If \(v\ge18\), then
\(3^v\ge3^{18}>2^{27}\), so \(x>27\). Ellison
\cite[Theorem~3]{ellison} gives
\[
  2^x-3^v>2^x e^{-x/10}>2^{x-v}.
\]
The last inequality follows from \(x/10\le v/5<v\log 2\).
\end{proof}

\begin{lemma}\label{lem:criterion}
Fix \(m\in\Npos\) and put \(n=1+\nu_2(3m-2)\). Write
\[
  3m-2=2^{n-1}u,\qquad u+1=2^v q,\qquad
  v=\nu_2(u+1).
\]
Here \(u,q\) are positive and odd and \(v\ge1\). Then \(S(m)\) is a
positive integer and
\begin{equation}\label{eq:criterion}
  S(m)>m
  \quad\Longleftrightarrow\quad
  3^v>2^{n+v}
  \quad\Longleftrightarrow\quad
  n+v>a(n).
\end{equation}
\end{lemma}

\begin{proof}
The identity \(3(4m-3)+1=4(3m-2)\) shows that \(F(4m-3)=u\).
Consequently
\[
  6S(m)=3+3^v q,\qquad 6m=2^{n+v}q-2^n+4.
\]
The first numerator is even and divisible by \(3\), so \(S(m)\)
is a positive integer. Subtraction gives
\begin{equation}\label{eq:difference}
  6(S(m)-m)=(3^v-2^{n+v})q+2^n-1.
\end{equation}
The two powers have opposite parity. If \(3^v>2^{n+v}\), the right-hand
side is positive. If \(3^v<2^{n+v}\), Lemma~\ref{lem:gap} and \(q\ge1\)
make it nonpositive. Finally,
\[
  3^v>2^{n+v}
  \quad\Longleftrightarrow\quad
  n+v>\frac{n\log 3}{\log(3/2)}.
\]
Since \(n+v\) is an integer, this is precisely \(n+v>a(n)\).
\end{proof}

\section{Proof of the bijection}

\begin{proof}[Proof of Theorem~\ref{thm:main}]
For a positive integer \(m\) and its associated row
\(n=1+\nu_2(3m-2)\), write \(a=a(n)\). The identity
\[
  3m-2+2^{n-1}=2^{n-1+v}q
\]
and the oddness of \(q\) show that the last condition in
\eqref{eq:criterion} is equivalent to \(2^a\) dividing
\(3m-2+2^{n-1}\). For an ascent we can therefore write
\begin{equation}\label{eq:t}
  3m=2^a t-2^{n-1}+2,\qquad t\in\Npos.
\end{equation}
Reduction modulo \(3\) gives
\[
 t\equiv
 \begin{cases}
  0\pmod{3}, & n\text{ even},\\
  (-1)^{a+1}\pmod{3}, & n\text{ odd}.
 \end{cases}
\]
Indeed, \(2^a t\equiv (-1)^{n-1}-2\pmod{3}\) and \(2^a\equiv
(-1)^a\pmod{3}\). Thus the positive solutions \(t\), in increasing
order, are
\begin{equation}\label{eq:column}
  t=3k-c,\qquad
  c=\frac{(3-(-1)^a)(1-(-1)^n)}{4},\qquad k\in\Npos.
\end{equation}
The correction \(c\) is \(0\) when \(n\) is even, \(1\) when \(n\) is odd
and \(a\) even, and \(2\) when both are odd. Substituting
\eqref{eq:column} into \eqref{eq:t} gives exactly \eqref{eq:array}.

Conversely, choose any \(n,k\in\Npos\), let \(t\) be given by
\eqref{eq:column}, and define \(m\) by \eqref{eq:t}. The congruence
makes \(m\) an integer, and \(a(n)\ge2n\) follows from \((3/2)^2<3\).
It follows that \(2^a t-2^{n-1}+2>0\). Moreover,
\[
  3m-2=2^{n-1}\bigl(2^{a-n+1}t-1\bigr).
\]
The factor in parentheses is positive and odd, so \(m\) recovers the
chosen row \(n\). Also \(3m-2+2^{n-1}=2^a t\), hence \(m\) is an ascent
by Lemma~\ref{lem:criterion}. For this \(m\), equation \eqref{eq:t}
fixes \(t\) and equation \eqref{eq:column} fixes \(k\).

The upward-antidiagonal reading assigns the index
\(T_{n+k-2}+k\) to \((n,k)\), where \(T_s=s(s+1)/2\).
For the diagonal with \(n+k-2=s\), the indices are
\(T_s+1,\ldots,T_{s+1}\). These intervals partition \(\Npos\),
so reading the array in this order lists each ascent once.
\end{proof}

\section{Lean formalization}\label{sec:lean}

The accompanying Lean~4 project \cite{lean} uses mathlib \cite{mathlib}.
It defines \(S\) over the rationals using \eqref{eq:step} and proves
that it agrees with a natural-number implementation. The signed
formula \eqref{eq:array} is expressed in the integers, and the
development proves that the displayed divisions are exact for
positive indices. The exponent \(a(n)\) is defined using the real
logarithm and floor.

The theorem \texttt{Jeffery.ascent\_iff\_occurs} identifies the image
of the array with the ascent set, and
\texttt{Jeffery.entry\_has\_unique\_position} establishes positivity
and injectivity. Theorems \texttt{Jeffery.jeffery\_conjecture} and
\texttt{Jeffery.jeffery\_sequence} give the unique-occurrence statements
in the array and in its upward-antidiagonal reading, respectively.
The proof of \texttt{Jeffery.dyadic\_gap} uses the argument in
Appendix~\ref{sec:formal-gap}. Merle's proof and its dependencies
\cite{merle} are included with attribution in the companion project.

The project was built with Lean~4.34.0 and mathlib~4.34.0. The kernel
also replayed the proof objects for all \(21\) project modules, including
the bundled logarithmic-bound modules. For the gap theorem, the
positivity and uniqueness theorem, and the three occurrence theorems
listed above, \texttt{\#print axioms} reports
\texttt{propext}, \texttt{Classical.choice}, and \texttt{Quot.sound}.
The companion archive \path{Jeffery-A254312-Lean.zip} contains the
source, pinned dependency manifest, build instructions, and
verification log.

\paragraph{AI assistance.}
OpenAI Codex assisted with proof development, Lean formalization,
and preparation of the manuscript.

\appendix
\section{The gap estimate used in Lean}\label{sec:formal-gap}

Assume \(n,v\ge1\) and \(3^v<2^x\), where \(x=n+v\). Put
\(\Delta=x\log 2-v\log 3>0\). If
\begin{equation}\label{eq:delta}
  \Delta\ge\frac{1}{2^v-1},
\end{equation}
then \(e^\Delta\ge1+\Delta\) yields
\(2^x/3^v\ge 2^v/(2^v-1)\). Rearranging gives the stronger inequality
\(2^x-3^v\ge 2^{x-v}\).

For \(v>20000\), Merle's theorem \texttt{Logs23.lam\_lower\_thm}
\cite{merle} bounds \(\lvert (v+n)\log 2-v\log 3\rvert\) below
by \(1/(2^v-1)\). The sign of \(\Delta\) gives \eqref{eq:delta}.

For \(18\le v\le20000\), the exact comparisons
\(2^{50508}<3^{31867}\) and \(3^{47468}<2^{75235}\) give
\[
  \frac{50508}{31867}<\frac{\log 3}{\log 2}
     <\frac{75235}{47468}.
\]
Suppose \(x/v\le75235/47468\). Since \(\Delta>0\), the integers
\(s=31867x-50508v\) and \(t=75235v-47468x\) satisfy \(s\ge1\)
and \(t\ge0\). The determinant identity
\(31867\cdot75235-50508\cdot47468=1\) gives
\(v=47468s+31867t\ge47468\), a contradiction. Therefore
\(47468x-75235v\ge1\), and the upper rational bound gives
\[
  \Delta>\frac{\log 2}{47468}
   >\frac{1}{94936}>\frac{1}{2^v-1}.
\]
Here \(\log 2=\int_1^2 dt/t>1/2\), and \(v\ge18\).
In the Lean development the two power comparisons are decided by
exact natural-number computation.

For \(1\le v\le17\), monotonicity in \(n\) reduces the check to the least
\(x\) with \(2^x>3^v\), as in Lemma~\ref{lem:gap}.
Table~\ref{tab:small} gives this \(x\) and the nonnegative excess
\(E=2^x-3^v-2^{x-v}+1\). In Lean, the same range is handled by
an elementary estimate for \(n>v\) and a kernel-checked finite
comparison for \(n\le v\). For the elementary estimate,
\(3^v\le4^v\le2^{n-1+v}\) and
\(2^{n-1+v}\ge2^n\) imply \(2^{n+v}-3^v\ge2^n\).

\begin{table}[htbp]
  \centering
  \begin{tabular}{@{}rrr@{\qquad}rrr@{}}
    \hline
    \(v\)&\(x\)&\(E\) & \(v\)&\(x\)&\(E\)\\
    \hline
    1&2&0 & 10&16&6\,424\\
    2&4&4 & 11&18&84\,870\\
    3&5&2 & 12&20&516\,880\\
    4&7&40 & 13&21&502\,574\\
    5&8&6 & 14&23&3\,605\,128\\
    6&10&280 & 15&24&2\,427\,798\\
    7&12&1\,878 & 16&26&24\,061\,120\\
    8&13&1\,600 & 17&27&5\,076\,542\\
    9&15&13\,022 & &&\\
    \hline
  \end{tabular}
  \caption{Exact excess at the least admissible exponent \(x\).}
  \label{tab:small}
\end{table}

\bibliographystyle{amsplain}
\bibliography{references}

\end{document}
